\documentclass[11pt]{article}
\usepackage[T1]{fontenc}
\usepackage[utf8]{inputenc}
\usepackage[a4paper,margin=25mm]{geometry}
\usepackage{amsmath,amssymb,amsthm,microtype,xcolor}
\usepackage[colorlinks=true,linkcolor=teal!70!black,urlcolor=teal!70!black,citecolor=teal!70!black]{hyperref}
\newtheorem{theorem}{Theorem}
\setlength{\parindent}{0pt}
\setlength{\parskip}{5pt}
\title{Time Translations and Finite Celestial Mellin Bands}
\hypersetup{pdftitle={Time Translations and Finite Celestial Mellin Bands}}
\author{\normalsize Deep Bhattacharjee, James Montgomery Lockwood, Tristan Romanov,\\
\normalsize Ikailo John Sessford, Ushashi Bhattacharya, Sanjeevan Singha Roy}
\date{Working research draft, 15 November 2026}
\begin{document}
\maketitle
\begin{center}\small\color{teal!70!black}
One unranked prospective author list supplied by the project lead.\\
Coauthor review and approval remain pending.
\end{center}
\begin{abstract}
We study exact frequency cutoffs on the celestial principal series
for a positive-energy massless scalar. Every nonzero time translation
sends each nonzero wave packet supported in a bounded Mellin band
outside every bounded output band. On any fixed band of positive
measure, the retained norm of a normalized packet can approach zero
or one. Every nonzero exactly band-limited packet also has infinite
mean physical energy. The statements concern one-particle kinematics.
\end{abstract}

\section{Introduction}
The four-dimensional conformal basis expresses the massless
one-particle energy variable by a Mellin transform \cite{basis}.
Translations act as multiplication by a phase in the energy
representation. Their formal action as shifts of the Mellin weight is
part of the Poincar\'e constraints on celestial amplitudes
\cite{lawzlotnikov}. The following result quantifies how much norm is
lost when one represents the state using only a fixed bounded band
of Mellin frequencies and then applies an ordinary time translation.

Both limits for retained norm come from shifting the same
log-energy packet. Such a shift preserves its Mellin support while
changing its physical energies. A separate analytic argument applies
to every nonzero packet and explains why exact hard frequency
cutoffs cannot give finite mean energy.

\section{Finite Mellin bands under time translation}
\label{sec:mellintranslation}
Consider a positive-energy massless scalar in four dimensions with
one-particle norm
\[
 \|f\|^2=\frac12\int_{S^2}\int_0^\infty |f(\omega,n)|^2\,
                   \omega\,d\omega\,dn .
\]
Take energies in a fixed unit, set $x=\log\omega$, and write
$g(x,n)=e^x f(e^x,n)/\sqrt 2$. The map $f\mapsto g$ is unitary to
$L^2(dx\,dn)$, and the unitary Fourier transform in $x$ gives the
Mellin representation on the principal line, with weight
$\Delta=1-i\nu$ for the sign convention $e^{-i\nu x}$.
This is the normalized wave-packet version of the massless conformal
basis \cite{basis}. Ordinary time translation acts by
\[
 (U_\tau g)(x,n)=e^{i\tau e^x}g(x,n).
\]
Here the sign of $\tau$ specifies our translation convention. The
factor $\omega$ in the generator also gives the familiar formal
unit shift of the Mellin weight \cite{lawzlotnikov}.

\begin{theorem}[Translation and finite-band retention]
\label{thm:mellintranslation}
Let $B\subset\mathbb R$ be bounded and measurable with positive
measure, and let $P_B$ be the orthogonal projector onto Mellin
frequencies in $B$. For each $\tau\ne0$,
\[
 \inf_{\substack{g=P_Bg\\ \|g\|=1}}\|P_BU_\tau g\|=0,
 \qquad
 \sup_{\substack{g=P_Bg\\ \|g\|=1}}\|P_BU_\tau g\|=1.
\]
For every nonzero $g$ in this band, $U_\tau g$ has nonzero norm
outside every bounded Mellin band. Also, for any $p>0$, every
nonzero $g$ in the band has
\[
 \int_{S^2}\int_{\mathbb R}e^{2px}|g(x,n)|^2\,dx\,dn=\infty.
\]
\end{theorem}
\begin{proof}
For the two extremal statements it suffices to choose a unit radial
profile $h$ with Fourier support in $B$ and a unit angular factor.
Write $h_R(x)=h(x-R)$, which preserves the Mellin support and norm.
After recentering $x$ by $R$, the retained norm for $U_\tau h_R$
is the norm of $\mathbf1_B\mathcal F(e^{i\tau e^R e^x}h)$.
For $R\to-\infty$, dominated convergence in $L^2(dx)$ gives
$\|(e^{i\tau e^R e^x}-1)h\|\to 0$, proving the supremum.
For $R\to+\infty$, first take an auxiliary profile $k\in C_c^\infty(\mathbb R)$.
Its transformed translated value is
\[
 \mathcal F(e^{i\tau e^R e^x}k)(\nu)
 =\frac1{\sqrt{2\pi}}\int_0^\infty
         e^{i\tau e^R y}k(\log y)y^{-i\nu-1}\,dy.
\]
The integrand with the oscillatory factor removed is in $L^1(dy)$.
The Riemann--Lebesgue lemma gives pointwise convergence to zero,
and the uniform bound $\|k\|_1/\sqrt{2\pi}$ together with
$|B|<\infty$ gives convergence in $L^2(B)$. Approximating $h$ by
such $k$ in $L^2$ extends the conclusion since the translation,
Fourier transform, and restriction have operator norms at most
one. This proves the infimum.

We use the following log-integrability fact for a nonzero
band-limited radial profile $h$. Write
\[
 h(z)=\frac1{\sqrt{2\pi}}\int_a^b \widehat h(\nu)e^{i\nu z}\,d\nu
\]
for some finite $a,b$. The integral defines an entire
representative since $\widehat h\in L^1([a,b])$.
For every real $y$, the function $F(z)=e^{-iaz}h(z+iy)$ is
bounded on the upper half-plane: in the integral its modulus
has the factor $e^{-(\nu-a)\Im z-\nu y}$. Choose
$z_0=x_0+iv_0$, $v_0>0$, with $F(z_0)\ne0$, and an upper
bound $M$ for $|F|$. Pull back to the unit disk by a
M\"obius map taking zero to $z_0$. The subharmonic mean-value
inequality for $\log|F|$ bounds the circular means of
$\log(M/|F|)$ by $\log(M/|F(z_0)|)$. Fatou's lemma,
continuity at every finite boundary point, and the Poisson
density for this map then give
\[
 \int_{\mathbb R}\log\frac{M}{|h(x+iy)|}\,
      \frac{v_0\,dx}{\pi((x-x_0)^2+v_0^2)}<\infty.
 \tag{*}\label{eq:bandlogintegrability}
\]
In particular, a bound $|h(x+iy)|\le C\exp(-ce^x)$
for all sufficiently large $x$, with $c>0$, is impossible.

If both $h$ and $U_\tau h$ belonged to bounded bands, both would
extend to entire functions, and their real-axis identity
would give $(U_\tau h)(z)=e^{i\tau e^z}h(z)$ everywhere.
Take $y=-\operatorname{sgn}(\tau)\pi/2$. On the line
$z=x+iy$ the entire representative of $U_\tau h$ is
bounded, because its compactly supported Fourier transform
is integrable. The multiplying factor has modulus
$\exp(|\tau|e^x)$ there, contrary to
\eqref{eq:bandlogintegrability}.

Finally, suppose $\int e^{2px}|h(x)|^2dx<\infty$. For real
$x$ set $u(x)=e^{px}|h(x)|$. On $x\ge0$,
\[
 \log\frac M{|h(x)|}\ \ge\ \log M+px-\log^+u(x).
\]
Since $\log^+u\le u^2$ and the Poisson density is bounded,
the last term is integrable with that weight. The term $px$
has divergent integral against the same density, again
contrary to \eqref{eq:bandlogintegrability}. For the two
contradictions with angular dependence, take a nonzero scalar
angular coefficient of $g$ in an orthonormal basis of $L^2(S^2)$.
Taking this coefficient commutes with $U_\tau$ and $P_B$, and
a finite weighted norm of $g$ would give a finite weighted
norm for each coefficient.
\end{proof}

For a unitary operator $U$ and an orthogonal projector $P$,
the exact norm deficit of its compression on $\operatorname{ran}P$ is
\[
 I_{\operatorname{ran}P}-(PUP)^*(PUP)
 =\big((I-P)UP\big)^*\big((I-P)UP\big).
\]
The two extremal cases above give operator norm one for both
the retained component $P_BU_\tau P_B$ and the discarded
component $(I-P_B)U_\tau P_B$. In energy terms, the moment
statement includes infinite mean energy for hard-banded
nonzero wave packets by taking $p=1/2$. These constraints
apply to an exact finite Mellin band in one-particle scalar
kinematics. Interacting infrared scattering and construction
of an independent boundary theory require separate results.

A numerical control can use the unit profile
$h_\lambda(x)=\sqrt{2\lambda}\exp(x/2-\lambda e^x)$, $\lambda>0$.
The gamma integral and $|\Gamma(\tfrac12+i\nu)|^2=\pi/\cosh(\pi\nu)$ give
\[
 |\mathcal F(U_\tau h_\lambda)(\nu)|^2
 = \frac{\lambda}{\sqrt{\lambda^2+\tau^2}}\,
   \frac{\exp(2\nu\arctan(\tau/\lambda))}{\cosh(\pi\nu)} .
\]
For $B=[-b,b]$ the initial squared norm in the band is
$q_b=2\arctan(\sinh(\pi b))/\pi$. The normalized band projection
of $h_\lambda$ differs from $h_\lambda$ in norm by
$\sqrt{2(1-\sqrt{q_b})}$. Contractivity of $P_BU_\tau$ then
bounds the projected packet's retained norm using the explicit
density, with the same error. For $b=8$ and $\tau=1$, the
numerical controls give a retained norm above $0.99999$ for
$\lambda=10^6$ and below $0.00006$ for $\lambda=10^{-10}$.


\section*{Reproducibility}
The repository includes checks of the finite-band example and a
quadrature of its Mellin amplitude. From the repository root they run
with \texttt{python scripts/verify\_mellin\_translation.py}.
\href{https://github.com/anaxia-research/celestial-infrared-stability}{Repository and source}.

\begin{thebibliography}{9}
\bibitem{basis} S. Pasterski and S.-H. Shao,
\emph{A conformal basis for flat space amplitudes},
Phys. Rev. D 96 (2017) 065022,
\href{https://arxiv.org/abs/1705.01027}{arXiv:1705.01027}.
\bibitem{lawzlotnikov} Y. T. A. Law and M. Zlotnikov,
\emph{Poincar\'e constraints on celestial amplitudes},
JHEP 03 (2020) 085,
\href{https://arxiv.org/abs/1910.04356}{arXiv:1910.04356}.
\end{thebibliography}
\end{document}
